% from ega4-9-fr.tex
\oldpage[IV-3]{73} $f_\eta(X_\eta)$ est localement fermé dans $Y_\eta$.
  Comme $f(X)$ est constructible dans $Y$ (1.8.4) et qu'il en est de même de
  $\overline{f(X)}$ puisque $Y$ est noethérien, il résulte de (8.3.11) qu'en
  restreignant au besoin $S$, on peut supposer que $f(X)$ est localement fermé
  dans $Y$. Il y a alors un ouvert $V$ de $Y$ contenant $f(X)$ et dans lequel
  $f(X)$ est fermé. Comme $Y$ est noethérien, $V$ est de type fini sur $S$, et
  en remplaçant $Y$ par $V$, on peut donc se ramener au cas où $f(X)$ est
  \emph{fermé} dans $Y$. Mais alors $f_s(X_s)=(f(X))_s$ est fermé dans $Y_s$
  pour tout $s\in S$, et dire que $f_s$ est une immersion équivaut à dire que
  $f_s$ est une immersion fermée~; on est donc ramené à ce qui a été démontré
  dans (ix).
\end{enumerate}

(x) Utilisant cette fois (2.7.1, (ix)) et (8.10.5, (iii)), on est ramené au cas
où $S$ est affine, noethérien, intègre de point générique $\eta$, et où
$\eta\in S-E$. Distinguons trois cas~:

\begin{enumerate}[label=\upshape{$\arabic{enumi}^\circ$}]
  \item $f_\eta(X_\eta)$ n'est pas ouvert dans $Y_\eta$. Comme $f(X)$ est
  constructible dans $Y$ (1.8.4), on déduit de (9.5.4) que $f_s(X_s)$ n'est
  pas ouvert dans $Y_s$ pour $s$ voisin de $\eta$, et \emph{a fortiori} $f_s$
  n'est pas une immersion ouverte.

  \item $f_\eta(X_\eta)$ est ouvert dans $Y_\eta$, mais $f_\eta$ n'est pas une
  immersion. Il résulte alors de (viii) que pour $s$ voisin de $\eta$ dans $S$,
  $f_s$ n'est pas une immersion, ni \emph{a fortiori} une immersion ouverte.

  \item $f_\eta(X_\eta)$ est ouvert dans $Y_\eta$ et $f_\eta$ est une
  immersion. Comme $f(X)$ est constructible dans $Y$, il résulte de (8.3.11)
  qu'en restreignant au besoin $S$, on peut déjà supposer que $f(X)$ est ouvert
  dans $Y$. Comme $Y$ est noethérien, le sous-préschéma induit sur $f(X)$ est
  de type fini sur $S$, donc on peut se ramener au cas où $f$ est \emph{surjectif}
  en remplaçant $Y$ par $f(X)$. Par hypothèse, $f_\eta$ est une immersion
  fermée, donc on peut, comme dans (ix), supposer que $f$ est une immersion
  \emph{fermée}, et par suite que $X$ est un sous-préschéma fermé de $Y$ défini
  par un Idéal cohérent $\mathcal J$ de $\mathcal O_Y$. Par hypothèse
  $f_\eta$ n'est pas un isomorphisme, donc $\mathcal J_\eta\neq0$~; on en
  conclut (9.4.5) que dans un voisinage de $\eta$, on a $\mathcal J_s\neq0$,
  et par suite l'immersion fermée surjective $f_s$ n'est pas ouverte.
\end{enumerate}

(xi) La propriété est conjonction des propriétés (i) et (x) et résulte donc de
ce qui a été démontré.

(xii) En vertu de (\textbf{I}, 5.3.8), dire que $f_s$ est un monomorphisme
signifie que
$\Delta_{f_s}=(\Delta_f)_s:X_s\to X_s\times_{Y_s}X_s=(X\times_YX)_s$ est un
isomorphisme, et comme $X\times_YX$ est un $S$-préschéma de présentation finie
(1.6.2, (iv)), la conclusion résulte de (xi).

\begin{proposition}[9.6.2]
\label{IV.9.6.2-fr}
Soient $X$, $Y$ deux $S$-préschémas de présentation finie, $f:X\to Y$ un
$S$-morphisme.
\begin{enumerate}[label=\upshape{I)}]
  \item Soit $E$ l'ensemble des $s\in S$ pour lesquels $f_s$ a l'une des
  propriétés suivantes~: être~:
  \begin{enumerate}[label=(\roman*)]
    \item affine~;
    \item quasi-affine~;
    \item projectif~;
    \item quasi-projectif.
  \end{enumerate}
  Alors $E$ est ind-constructible dans $S$.
\end{enumerate}

